\chapter{Ejercicios de Sistemas Distribuidos}

% Configuración de estilos para el código
\definecolor{backcolour}{rgb}{0.95,0.95,0.92}
\definecolor{codegray}{rgb}{0.5,0.5,0.5}
\definecolor{keyblue}{rgb}{0.1,0.1,0.8}

\lstset{
    backgroundcolor=\color{backcolour},   
    basicstyle=\ttfamily\small,
    breakatwhitespace=false,         
    breaklines=true,                 
    captionpos=b,                    
    commentstyle=\color{gray},
    keywordstyle=\color{keyblue},
    numbers=left,                    
    numberstyle=\tiny\color{codegray},
    stepnumber=1,
    numbersep=8pt,                  
    showspaces=false,                
    showstringspaces=false,
    showtabs=false,                  
    tabsize=2,
    frame=single, 
    rulecolor=\color{black!30}
}

% --- SECCIÓN 1: ASIGNADOR DE RECURSOS ---
\section{Asignador de Recursos}
Diseño de un asignador de recursos con gestión de cola de espera manual (Paso de Mensajes).


\begin{lstlisting}[language=Pascal, caption=Asignador de Recursos]
// Definicion de canales
type clase_op = enum (ADQUIRIR, LIBERAR);
port peticion (int id_cliente, clase_op operacion, int unidad);
chan respuesta [1..n](int unidad_asignada);

Servidor {
    var unidades: set of int; 
    var pendientes: queue of int;
    var disponibles: int := MAX_UNIDADES;

    while true {
        // 1. PRIORIDAD: Si hay alguien esperando y hay recursos, desbloquear
        if (not empty(pendientes) && disponibles > 0) {
            cliente = dequeue(pendientes);
            unidad = extraer(unidades);
            disponibles--;
            send respuesta[cliente](unidad);
        }

        // 2. ESCUCHA: El servidor atiende peticiones
        receive peticion(i, operacion, uni)
        switch operacion {
            case ADQUIRIR:
                if (disponibles > 0) {
                    send respuesta[i](extraer(unidades));
                    disponibles--;
                } else {
                    // Bloqueo logico: No hay send, el cliente espera
                    enqueue(pendientes, i);
                }
            case LIBERAR:
                if (empty(pendientes)) {
                    insertar(unidades, uni);
                    disponibles++;
                } else {
                    // Entrega directa al siguiente en la cola
                    proximo = dequeue(pendientes);
                    send respuesta[proximo](uni);
                }
        }
    }
}
\end{lstlisting}

% --- SECCIÓN 2: SERVIDOR DE FICHEROS ---
\section{Servidor de Ficheros}
Gestión de sesiones mediante canales dinámicos y mailbox.

\begin{lstlisting}[caption=Servidor de Ficheros]
// Canales: Mailbox para apertura, array para operaciones
mailbox abrir(string nombref, int id_cli);
chan leer[1..n](...), escribir[1..n](...), cerrar[1..n](...);
chan respuesta_abrir[1..n](int descriptor);
do receive abrir(nombref, id_cli) ->
    // Se asigna el canal 'i' (descriptor de fichero)
    send respuesta_abrir[id_cli](i); 
    
    do // Bucle de sesion privada para el fichero i
        receive leer[i](datos) -> 
            read(i, datos); 
            send respuesta_cli[id_cli](datos);
            
        [] receive escribir[i](datos) -> 
            write(i, datos);
            send respuesta_cli[id_cli](ok);
            
        [] receive cerrar[i]() -> 
            close(i);
            send respuesta_cli[id_cli](bye);
            break; // Salida del bucle interno
    od
    fichero_abierto[i] = false;
od
\end{lstlisting}

% --- SECCIÓN 3: BUFFER ACOTADO ---
\section{Buffer Acotado}
Implementación mediante Órdenes Guardadas (Estilo Ada/CSP).

\begin{lstlisting}[caption=Buffer Acotado con Guardas]
Bufferacotado::
    op poner(dato: T);
    op tomar(var res: T);
    
    var buffer[1..n]: T;
    var cab: int := 1; cola: int := 1; contador: int := 0;

    do true ->
        in poner(dato) and (contador < n) ->
            buffer[cola] := dato;
            cola := (cola mod n) + 1;
            contador++;
        [] tomar(res) and (contador > 0) ->
            res := buffer[cab];
            cab := (cab mod n) + 1;
            contador--;
        nin
    od
\end{lstlisting}

% --- SECCIÓN 4: BANCO (CON IF) ---
\section{Ejercicio Examen Banco}
Servidor con estado, tokens de seguridad y lógica de control mediante IF.

mailbox identificar(string nombre, string contra, int j);
// Operaciones privadas: Un array de canales, uno para cada servidor [1..n]
chan SeleccionarCuenta[1..n](int token, TCuenta cuenta, int j);
chan ConsultarSaldo[1..n](int token, int j);
chan IngresarDinero[1..n](int token, float cantidad, int j);
chan RetirarDinero[1..n](int token, float cantidad, int j);
chan BloquearCuenta[1..n](int token, int j);
chan SalirCliente[1..n](int token, int j);

// Respuestas al cliente: Un array de canales para los 'm' clientes del sistema
chan respuesta_id[1..m](int token, TLista lista);
chan respuesta_op[1..m](any resultado);
\begin{lstlisting}[caption=Servidor Bancario con Estado]
Servidor::[i:1..n]
{
    // Variables locales del proceso i
    var salir[i]: boolean := false;
    var seleccionada[i]: TCuenta := 0;
    var tipo[i]: int := 0;
    var token_sesion[i]: int;

    while (true)
    {
        // Fase 1: Identificación (Cualquier servidor libre atiende)
        receive identificar(nombre, contrasena, j); 

        if (comprobar(nombre, contrasena)) 
        {
            send respuesta_id[j](token_sesion[i], listaCuentas);
            salir[i] := false; 

            // Fase 2: Sesión Privada (Bucle de operaciones)
            while(salir[i] == false) 
            {
                // Operación sin guarda: siempre se puede seleccionar cuenta
                if receive SeleccionarCuenta[i](token, cuenta, j) ->
                    seleccionada[i] = cuenta;
                    tipo[i] = cuenta.TipoAcceso;
                    send respuesta_op[j](tipo[i]);

                // Operaciones con GUARDAS: Solo si hay cuenta seleccionada != 0
                
                // Consultar Saldo
                [] seleccionada[i] != 0; receive ConsultarSaldo[i](token, j) ->
                    send respuesta_op[j](seleccionada[i].saldo);

                // Ingresar Dinero (Requiere cuenta y permiso de escritura 'tipo')
                [] seleccionada[i] != 0 && tipo[i]; receive IngresarDinero[i](token, cantidad, j) ->
                    seleccionada[i].saldo += cantidad;
                    send respuesta_op[j](seleccionada[i].saldo);

                // Retirar Dinero (Uso de IF interno para control de saldo)
                [] seleccionada[i] != 0 && tipo[i]; receive RetirarDinero[i](token, cantidad, j) ->
                    if (cantidad <= seleccionada[i].saldo) {
                        seleccionada[i].saldo -= cantidad;
                        send respuesta_op[j](seleccionada[i].saldo);
                    } else {
                        send respuesta_op[j](0); // Error: Saldo insuficiente
                    }

                // Bloquear Cuenta
                [] seleccionada[i] != 0 && tipo[i]; receive BloquearCuenta[i](token, j) ->
                    seleccionada[i] = 0;
                    send respuesta_op[j](resultado);

                // Salir Cliente: Cierra el bucle interno y limpia estado
                [] seleccionada[i] != 0 && tipo[i]; receive SalirCliente[i](token, j) ->
                    salir[i] = true;
                    send respuesta_op[j](ok);
                    seleccionada[i] = 0;
                    tipo[i] = 0;
            }
        }
    }
}
\end{lstlisting}